\documentclass[UTF8,a4paper,12pt]{ctexart}

% 加载报告排版宏包
\usepackage[left=2.6cm,right=2.6cm,top=2.6cm,bottom=2.6cm,headheight=15pt]{geometry}
\usepackage{setspace}
\usepackage{graphicx}
\usepackage{subcaption}
\usepackage{float}
\usepackage{booktabs}
\usepackage{array}
\usepackage{tabularx}
\usepackage{longtable}
\usepackage{multirow}
\usepackage{enumitem}
\usepackage{chngcntr}
\usepackage{xcolor}
\usepackage{tikz}
\usepackage{caption}
\usepackage{fancyhdr}
\usepackage{tocloft}
\usepackage{hyperref}
\usepackage{listings}

% 加载绘图组件
\usetikzlibrary{arrows.meta,positioning,calc,fit,shapes.geometric,backgrounds,decorations.pathreplacing}

% 设置图片路径
\graphicspath{{figure/}}

% 设置超链接样式
\hypersetup{
  colorlinks=true,
  linkcolor=black,
  urlcolor=black,
  citecolor=black
}

% 设置正文格式
\onehalfspacing
\zihao{-4}
\setlength{\parindent}{2em}
\setlength{\parskip}{0pt}
\setlist{nosep,leftmargin=2.5em}

% 设置章节标题格式
\ctexset{
  section={
    format={\centering\heiti\zihao{3}},
    name={,},
    number=\arabic{section},
    aftername=\quad,
    beforeskip=1.2em,
    afterskip=1.0em
  },
  subsection={
    format={\heiti\zihao{4}},
    name={,},
    number=\arabic{section}.\arabic{subsection},
    aftername=\quad,
    beforeskip=0.8em,
    afterskip=0.4em
  },
  subsubsection={
    format={\heiti\zihao{-4}},
    name={,},
    number=\arabic{section}.\arabic{subsection}.\arabic{subsubsection},
    aftername=\quad,
    beforeskip=0.5em,
    afterskip=0.2em
  }
}

% 设置目录格式
\renewcommand{\contentsname}{\zihao{3}\heiti 目\quad 录}
\renewcommand{\cfttoctitlefont}{\hfill\zihao{3}\heiti}
\renewcommand{\cftaftertoctitle}{\hfill}
\renewcommand{\cftsecfont}{\songti\zihao{-4}}
\renewcommand{\cftsubsecfont}{\songti\zihao{-4}}
\renewcommand{\cftsubsubsecfont}{\songti\zihao{-4}}
\renewcommand{\cftsecpagefont}{\songti\zihao{-4}\mdseries}
\renewcommand{\cftsubsecpagefont}{\songti\zihao{-4}\mdseries}
\renewcommand{\cftsubsubsecpagefont}{\songti\zihao{-4}\mdseries}
\renewcommand{\cftsecleader}{\cftdotfill{\cftdotsep}}
\setlength{\cftbeforetoctitleskip}{0pt}
\setlength{\cftaftertoctitleskip}{0.8em}
\setlength{\cftbeforesecskip}{0.5em}
\setlength{\cftbeforesubsecskip}{0.18em}
\setlength{\cftbeforesubsubsecskip}{0.12em}
\cftsetindents{section}{0em}{2.2em}
\cftsetindents{subsection}{2.0em}{3.0em}
\cftsetindents{subsubsection}{4.4em}{4.2em}

% 设置图表格式
\renewcommand{\figurename}{图}
\renewcommand{\tablename}{表}
\captionsetup{font=small,labelfont=bf,labelsep=quad,justification=centering}
\counterwithin{figure}{section}
\counterwithin{table}{section}
\renewcommand{\thefigure}{\thesection-\arabic{figure}}
\renewcommand{\thetable}{\thesection-\arabic{table}}

% 设置代码摘录格式
\lstset{
  basicstyle=\ttfamily\small,
  keywordstyle=\color{blue}\bfseries,
  commentstyle=\color{gray},
  stringstyle=\color{red},
  numbers=left,
  numberstyle=\tiny\color{gray},
  stepnumber=1,
  numbersep=5pt,
  backgroundcolor=\color{white},
  showspaces=false,
  showstringspaces=false,
  showtabs=false,
  frame=single,
  rulecolor=\color{black},
  tabsize=4,
  captionpos=b,
  breaklines=true,
  breakatwhitespace=false,
  columns=fullflexible,
  keepspaces=true,
  escapeinside={\%*}{*)},
  xleftmargin=2em,
  xrightmargin=2em,
  aboveskip=1em,
  belowskip=1em
}

\lstdefinestyle{VerilogStyle}{
  language=Verilog,
  keywordstyle=\color{blue}\bfseries,
  commentstyle=\color{green!60!black},
  stringstyle=\color{orange},
  breaklines=true,
  breakatwhitespace=false
}

\lstdefinestyle{CSharpStyle}{
  language=[Sharp]C,
  keywordstyle=\color{blue}\bfseries,
  commentstyle=\color{green!60!black},
  stringstyle=\color{orange},
  breaklines=true,
  breakatwhitespace=false
}

% 设置页眉页脚
\pagestyle{fancy}
\fancyhf{}
\fancyhead[C]{\songti\zihao{-5}基于FPGA的HDB3编译码器系统设计}
\fancyfoot[C]{\songti\zihao{-5}\thepage}
\renewcommand{\headrulewidth}{0.4pt}
\renewcommand{\footrulewidth}{0pt}

\fancypagestyle{plain}{
  \fancyhf{}
  \fancyhead[C]{\songti\zihao{-5}基于FPGA的HDB3编译码器系统设计}
  \fancyfoot[C]{\songti\zihao{-5}\thepage}
  \renewcommand{\headrulewidth}{0.4pt}
  \renewcommand{\footrulewidth}{0pt}
}

% 定义图形节点样式
\tikzset{
  block/.style={draw=blue!55!black,rounded corners=2pt,align=center,minimum height=0.9cm,minimum width=2.5cm,fill=blue!5,line width=0.6pt},
  smallblock/.style={draw=blue!55!black,rounded corners=2pt,align=center,minimum height=0.75cm,minimum width=1.8cm,fill=blue!5,font=\small,line width=0.6pt},
  ioblock/.style={draw=teal!60!black,rounded corners=2pt,align=center,minimum height=0.9cm,minimum width=2.5cm,fill=teal!6,line width=0.6pt},
  coreblock/.style={draw=orange!70!black,rounded corners=2pt,align=center,minimum height=0.9cm,minimum width=2.5cm,fill=orange!8,line width=0.6pt},
  storeblock/.style={draw=violet!70!black,rounded corners=2pt,align=center,minimum height=0.8cm,minimum width=2.4cm,fill=violet!6,line width=0.6pt},
  decision/.style={diamond,aspect=2.1,draw=orange!70!black,align=center,inner sep=1pt,minimum height=0.9cm,fill=orange!8,line width=0.6pt},
  framebyte/.style={draw=blue!55!black,minimum height=0.72cm,align=center,fill=blue!5,line width=0.5pt},
  arrow/.style={-{Stealth[length=2.4mm]},thick,draw=black!75},
  bus/.style={-{Stealth[length=2.4mm]},thick,draw=blue!65!black},
  returnbus/.style={-{Stealth[length=2.4mm]},thick,draw=teal!65!black},
  wavegrid/.style={gray!35,line width=0.35pt}
}

\begin{document}

% 生成目录页
\newpage
\pagenumbering{Roman}

\fancypagestyle{tocpage}{
  \fancyhf{}
  \fancyhead[C]{\songti\zihao{-5}基于FPGA的HDB3编译码器系统设计}
  \fancyfoot[C]{\songti\zihao{-5}\thepage}
  \renewcommand{\headrulewidth}{0.4pt}
  \renewcommand{\footrulewidth}{0pt}
}
\thispagestyle{tocpage}

\tableofcontents

\newpage
\pagenumbering{arabic}
\setcounter{page}{1}

\phantomsection
\section*{摘要}
\addcontentsline{toc}{section}{摘要}

本文围绕“HDB3 编译码器的仿真”课题，设计并实现了一套基于 FPGA 的 HDB3 编译码器系统。系统以单极性不归零二进制数字信号为输入，在 FPGA 中完成 HDB3 编码、译码、串口通信和双通道 DAC 波形输出；上位机采用 C\# WPF 编写，用于输入测试序列、计算软件参考结果、发送串口命令、接收 FPGA 应答并进行逐项比对。硬件平台采用 ACX720 开发板，FPGA 芯片为 GW5AT-138B，通信接口为 115200 bps、8N1 单 UART。

系统设计中，HDB3 编码器按照 AMI 极性交替和四连零替换规则生成 \texttt{+1}、\texttt{-1}、\texttt{+V}、\texttt{-V}、\texttt{+B}、\texttt{-B} 等符号；译码器将普通正负脉冲恢复为 1，将零电平、破坏脉冲和补信脉冲恢复为 0。FPGA 通过 \texttt{55 AA} 命令帧接收上位机请求，通过 \texttt{AA 55} 应答帧返回处理结果。双通道 DAC 中，DA0 输出 HDB3 编码波形（双极性 RZ），DA1 输出译码恢复的二进制波形（单极性 NRZ），便于示波器对比观察编码前后的信号变化。

实验部分完成了上位机与 FPGA 编码测试、上位机与 FPGA 译码测试以及示波器实测验证，并以 Vivado 仿真工程的 RTL 框图和 ModelSim 波形作为编译码逻辑的补充验证。全 0、全 1、10 交替、典型混合、长 0、编码随机和解码随机等测试结果表明，FPGA 返回结果与上位机软件参考结果一致，示波器波形能够体现 HDB3 码的正负交替、四连零替换和译码恢复特性。实验验证了本系统 HDB3 编译码功能的正确性。

\vspace{0.6em}
\noindent\textbf{关键词：}HDB3；FPGA；UART；DAC；上位机；示波器

\section{绪论}

\subsection{课题背景}

在数字基带传输系统中，线路编码用于把二进制数据转换为适合传输和观察的电平波形。不同编码方式会影响信号的直流分量、定时信息和波形可识别性。单极性不归零码实现简单，但当输入数据中出现较长连续 0 时，线路电平会长时间保持不变，接收端不容易从波形中提取同步信息。

AMI 码将二进制 1 交替表示为正、负脉冲，可以减小直流分量，但连续 0 仍然会造成长时间无跳变。HDB3 码在 AMI 码基础上加入四连零替换规则，通过破坏脉冲和补信脉冲限制连续零电平的长度。因此，HDB3 码既能保持 AMI 码正负脉冲交替的特点，又能改善长连零条件下的定时信息不足问题。

本设计围绕 HDB3 编码、译码和波形验证展开。为了让编码规则不只停留在理论推导中，本文将 HDB3 编译码逻辑放入 FPGA 实现，并配套设计上位机通信和 DAC 波形输出，使软件计算结果、FPGA 返回结果和示波器实测波形能够相互对应。

\newpage
\subsection{设计任务}

本设计以单极性不归零二进制序列为输入，在 FPGA 中完成 HDB3 编码和译码，并通过上位机与示波器完成结果验证。系统需要完成以下内容：

\begin{enumerate}
  \item 实现二进制序列到 HDB3 符号序列的编码。
  \item 实现 HDB3 符号序列到二进制序列的译码。
  \item 建立上位机与 FPGA 之间的单 UART 通信链路。
  \item 在上位机中生成软件参考结果，并与 FPGA 返回结果进行比对。
  \item 通过双通道 DAC 输出 HDB3 编码波形和译码恢复波形。
  \item 结合仿真、串口联调和示波器实测验证系统功能。
\end{enumerate}

\subsection{系统验证思路}

本文采用“软件参考、FPGA 返回、示波器波形”三层验证方式。首先，上位机软件根据 HDB3 编译码规则计算期望结果；然后，FPGA 对相同输入进行处理，并通过 UART 返回实际结果，上位机逐 bit 或逐符号进行比对；最后，FPGA 将编码结果和译码结果送入 DAC，通过示波器观察 DA0、DA1 两路波形是否符合预期。

此外，报告保留独立仿真工程结果，用于补充说明 HDB3 编译码逻辑在仿真环境中的结构和基本功能。该仿真工程不包含最终系统中的上位机通信、UART 协议解析和 DAC 播放结构，因此本文将其作为功能补充验证，而不作为最终 FPGA 系统结构的依据。

\section{HDB3 编译码原理}

\subsection{AMI 码原理}

AMI 码的编码规则为：输入 0 时输出零电平，输入 1 时输出非零脉冲，并且相邻的 1 依次采用正、负极性交替表示。若输入序列为连续的 1，则输出波形会呈现正、负、正、负交替；若输入序列中出现连续 0，则输出保持零电平。

AMI 码能够减少直流分量，但它对连续 0 没有额外处理。当连续 0 的长度较长时，传输线上长时间没有电平跳变，不利于接收端保持时钟同步。HDB3 码通过限制连续零电平个数，解决了这一问题。

\subsection{HDB3 编码规则}

HDB3 码的核心规则是在 AMI 码基础上对四连零进行替换，使输出序列中最多只出现三个连续零电平。工程实现中常采用奇偶法判断替换形式：自上一次破坏脉冲 \texttt{V} 以后，如果普通 1 的个数为偶数，则连续四个 0 替换为 \texttt{B00V}；如果普通 1 的个数为奇数，则连续四个 0 替换为 \texttt{000V}。

\begin{figure}[H]
  \centering
  \begin{tikzpicture}[font=\small,node distance=1.0cm]
    \node[smallblock,minimum width=2.6cm] (zeros) {检测到\\\texttt{0000}};
    \node[decision,right=1.1cm of zeros,minimum width=2.5cm] (judge) {普通 1\\个数};
    \node[coreblock,right=1.3cm of judge,yshift=0.88cm,minimum width=3.1cm] (even) {偶数\\\texttt{0000 -> B00V}};
    \node[coreblock,right=1.3cm of judge,yshift=-0.88cm,minimum width=3.1cm] (odd) {奇数\\\texttt{0000 -> 000V}};
    \node[smallblock,right=1.2cm of even,minimum width=2.8cm] (even_note) {插入补信脉冲\\保持极性平衡};
    \node[smallblock,right=1.2cm of odd,minimum width=2.8cm] (odd_note) {插入破坏脉冲\\标记替换位置};
    \draw[arrow] (zeros) -- (judge);
    \draw[arrow] (judge.east) -- ++(0.35,0) |- node[pos=0.24,above]{偶数} (even.west);
    \draw[arrow] (judge.east) -- ++(0.35,0) |- node[pos=0.24,below]{奇数} (odd.west);
    \draw[arrow] (even) -- (even_note);
    \draw[arrow] (odd) -- (odd_note);
  \end{tikzpicture}
  \caption{HDB3 四连零替换规则}
  \label{fig:hdb3_rule}
\end{figure}

其中，\texttt{V} 是破坏脉冲，其极性故意破坏 AMI 极性交替规律，使接收端能够识别四连零替换位置；\texttt{B} 是补信脉冲，用于维持整体直流平衡。替换完成后，编码器清零连续零计数和奇偶计数，重新开始统计。

需要说明的是，不同教材对 HDB3 编码规则的叙述方式有所差异。部分教材规定每次遇到四连零时应先判断该组替换是否为首个替换组———若是，则固定采用 \texttt{000V}；若非首个，再根据两次替换组之间的普通 1 个数奇偶性决定。本设计采用的奇偶法不区分首次与后续，统一按普通 1 的奇偶个数判断，初始奇偶状态为偶数。两种写法在绝大多数输入序列下输出一致，仅在个别边界条件下（如全零序列的首个替换组）替换形式有所不同。由于本系统的软件参考和 FPGA 硬件均以同一套奇偶法为准，软硬件对比结果始终保持一致，不影响系统功能的正确性验证。

\subsection{HDB3 译码规则}

HDB3 译码的目标是恢复原始二进制序列。普通正脉冲 \texttt{+1} 和普通负脉冲 \texttt{-1} 来源于原始输入中的 1，因此译码为 1；零电平 \texttt{0}、破坏脉冲 \texttt{+V/-V} 和补信脉冲 \texttt{+B/-B} 均对应原始输入中的 0，因此译码为 0。

本设计的译码器采用直接符号映射方式，不再反向查找 \texttt{B00V} 或 \texttt{000V} 的组合位置。该方式与本系统采用的编码符号定义一致，便于 FPGA 批处理实现和上位机逐项比较。

\subsection{符号定义}

上位机与 FPGA 通信时，每个 HDB3 符号使用一个字节表示，具体定义如表 \ref{tab:symbol_definition} 所示。

\begin{table}[H]
  \centering
  \caption{HDB3 符号定义}
  \label{tab:symbol_definition}
  \begin{tabular}{ccc}
    \toprule
    字节值 & 显示符号 & 含义 \\
    \midrule
    \texttt{0x00} & \texttt{0} & 零电平 \\
    \texttt{0x01} & \texttt{+1} & 正普通脉冲 \\
    \texttt{0x02} & \texttt{-1} & 负普通脉冲 \\
    \texttt{0x03} & \texttt{+V} & 正破坏脉冲 \\
    \texttt{0x04} & \texttt{-V} & 负破坏脉冲 \\
    \texttt{0x05} & \texttt{+B} & 正补信脉冲 \\
    \texttt{0x06} & \texttt{-B} & 负补信脉冲 \\
    \bottomrule
  \end{tabular}
\end{table}

\section{系统方案设计}

\subsection{总体结构}

系统由上位机、UART 通信链路、FPGA 下位机、双通道 DAC 和示波器组成。上位机输入测试序列并计算软件参考结果，FPGA 接收命令后完成编码或译码，并同时返回数字结果与输出模拟波形。总体结构如图 \ref{fig:system_structure} 所示。

\begin{figure}[H]
  \centering
  \resizebox{0.86\textwidth}{!}{%
  \begin{tikzpicture}[font=\small,node distance=1.05cm]
    \node[ioblock,minimum width=3.0cm,minimum height=1.05cm] (pc) {上位机\\输入与比对};
    \node[ioblock,right=of pc,minimum width=2.4cm,minimum height=1.05cm] (uart) {单 UART\\115200 bps};
    \node[storeblock,right=of uart,minimum width=3.0cm,minimum height=1.05cm] (fpga) {FPGA\\HDB3 编译码};
    \node[coreblock,right=of fpga,minimum width=2.4cm,minimum height=1.05cm] (dac) {双通道 DAC};
    \node[ioblock,right=of dac,minimum width=2.4cm,minimum height=1.05cm] (scope) {示波器};

    \draw[bus] (pc) -- node[above]{命令} (uart);
    \draw[bus] (uart) -- (fpga);
    \draw[bus] (fpga) -- node[above]{波形} (dac);
    \draw[bus] (dac) -- (scope);
    \draw[returnbus] ([yshift=0.16cm]fpga.north) -- ++(0,0.55) -| node[pos=0.65,above]{应答} ([yshift=0.16cm]pc.north);
  \end{tikzpicture}
  }
  \caption{系统总体结构}
  \label{fig:system_structure}
\end{figure}

系统实现时以 ACX720 开发板为硬件平台，FPGA 工程运行在 50 MHz 系统时钟下。串口通信采用 115200 bps、8N1 格式，DAC 播放模块将一个码元划分为两个采样阶段，因此 HDB3 码元速率为 100 ksym/s，DAC 采样时钟为 200 kHz。主要参数如表 \ref{tab:system_params} 所示。

\begin{table}[H]
  \centering
  \caption{系统主要设计参数}
  \label{tab:system_params}
  \begin{tabularx}{\textwidth}{c c X}
    \toprule
    参数项 & 取值 & 说明 \\
    \midrule
    硬件平台 & ACX720 / GW5AT-138B & FPGA 下位机实现 HDB3 编码、译码和 DAC 播放 \\
    系统时钟 & 50 MHz & 顶层逻辑、串口模块和 DAC 播放控制共用时钟 \\
    串口格式 & 115200 bps，8N1 & PC 与 FPGA 使用单 UART 进行命令和应答传输 \\
    命令帧头 & \texttt{55 AA} & PC 到 FPGA 的命令帧起始标志 \\
    应答帧头 & \texttt{AA 55} & FPGA 到 PC 的应答帧起始标志 \\
    序列长度 & 最多 64 个 bit / 符号 & 与 FPGA 缓存和应答帧长度一致 \\
    码元速率 & 100 ksym/s & 由 50 MHz 时钟经 500 个周期分频得到 \\
    DAC 采样率 & 200 kHz & 每个码元包含前半码元和后半码元两次采样 \\
    DAC 通道 & DA0 / DA1 & DA0 输出 HDB3 编码波形，DA1 输出译码恢复的二进制波形 \\
    \bottomrule
  \end{tabularx}
\end{table}

该结构将数字比对和波形观察结合起来。上位机比对能够精确判断每一个符号或 bit 是否正确；示波器波形能够直观展示 HDB3 码对长连零的处理效果和译码恢复结果。图 \ref{fig:hardware_connection} 为系统实物连接关系，ACX720 开发板通过 USB 转串口与 PC 通信，双通道 DAC 经 BNC 线接入示波器的 CH1、CH2，分别观察 DA0 和 DA1 的输出波形。

\begin{figure}[H]
  \centering
  \includegraphics[width=0.7\textwidth]{实物连线图.png}
  \caption{系统实物连接关系}
  \label{fig:hardware_connection}
\end{figure}

\subsection{FPGA 端结构}

FPGA 端不是简单的单模块编码器，而是由串口收发、命令解析、编译码核心、应答发送和 DAC 播放控制共同组成。最终 Gowin 工程的顶层模块为 \texttt{hdb3\_top}，模块关系如图 \ref{fig:fpga_structure} 所示。

\begin{figure}[H]
  \centering
  \resizebox{0.84\textwidth}{!}{%
  \begin{tikzpicture}[font=\small,node distance=0.9cm]
    \node[ioblock,minimum width=2.2cm] (rx) {UART 接收};
    \node[smallblock,right=of rx,minimum width=2.7cm] (parser) {命令解析};
    \node[coreblock,right=of parser,minimum width=2.9cm] (codec) {HDB3\\编码 / 译码};
    \node[storeblock,right=of codec,minimum width=2.7cm] (result) {结果缓存};
    \node[ioblock,right=of result,yshift=0.62cm,minimum width=2.3cm] (tx) {UART 应答};
    \node[ioblock,right=of result,yshift=-0.62cm,minimum width=2.3cm] (dacout) {DA0 / DA1};

    \draw[bus] (rx) -- (parser);
    \draw[bus] (parser) -- (codec);
    \draw[bus] (codec) -- (result);
    \draw[returnbus] ([yshift=0.14cm]result.east) -- ++(0.35,0) |- (tx.west);
    \draw[bus] ([yshift=-0.14cm]result.east) -- ++(0.35,0) |- (dacout.west);
  \end{tikzpicture}
  }
  \caption{最终 FPGA 工程模块关系}
  \label{fig:fpga_structure}
\end{figure}

\begin{table}[H]
  \centering
  \caption{FPGA 主要模块说明}
  \label{tab:fpga_modules}
  \begin{tabularx}{\textwidth}{>{\ttfamily}l>{\ttfamily}lX}
    \toprule
    模块 & 文件 & 功能 \\
    \midrule
    hdb3\_top & src/hdb3\_top.v & 顶层调度、缓存管理、状态机控制和子模块例化 \\
    packet\_parser & src/hdb3/packet\_parser.v & 解析 PC 到 FPGA 的 \texttt{55 AA} 命令帧 \\
    response\_tx & src/hdb3/response\_tx.v & 封装并发送 FPGA 到 PC 的 \texttt{AA 55} 应答帧 \\
    hdb3\_encoder & src/hdb3/hdb3\_encoder.v & 将二进制 bit 编码为 HDB3 符号 \\
    hdb3\_decoder & src/hdb3/hdb3\_decoder.v & 将 HDB3 符号译码为二进制 bit \\
    dac\_playback & src/hdb3/dac\_playback.v & 控制 DA0、DA1 双通道循环播放 \\
    uart\_byte\_rx & src/comm/uart\_byte\_rx.v & UART 字节接收 \\
    uart\_byte\_tx & src/comm/uart\_byte\_tx.v & UART 字节发送 \\
    \bottomrule
  \end{tabularx}
\end{table}

\subsection{UART 协议设计}

系统采用单 UART 通信，命令和应答共用同一串口。为了区分通信方向，PC 到 FPGA 的命令帧以 \texttt{55 AA} 开头，FPGA 到 PC 的应答帧以 \texttt{AA 55} 开头。帧格式如图 \ref{fig:uart_frame} 所示。

\begin{figure}[H]
  \centering
  \begin{tikzpicture}[font=\small]
    \node[anchor=east] at (-0.25,0.36) {命令帧};
    \foreach \name/\w/\text/\color [count=\i] in {
      head_a/1.05/\texttt{55}/teal!10,
      head_b/1.05/\texttt{AA}/teal!10,
      cmd/1.15/\texttt{cmd}/blue!5,
      len_l/1.2/\texttt{len\_l}/blue!5,
      len_h/1.2/\texttt{len\_h}/blue!5,
      payload/2.15/\texttt{payload}/orange!10,
      checksum/1.55/\texttt{checksum}/violet!8
    }{
      \ifnum\i=1
        \node[framebyte,minimum width=\w cm,fill=\color,anchor=west] (\name) at (0,0) {\text};
      \else
        \node[framebyte,minimum width=\w cm,fill=\color,anchor=west] (\name) at (prev.east) {\text};
      \fi
      \coordinate (prev) at (\name.east);
    }

    \node[anchor=east] at (-0.25,-1.04) {应答帧};
    \foreach \name/\w/\text/\color [count=\i] in {
      resp_head_a/1.05/\texttt{AA}/teal!10,
      resp_head_b/1.05/\texttt{55}/teal!10,
      resp_cmd/1.15/\texttt{cmd}/blue!5,
      status/1.25/\texttt{status}/red!8,
      resp_len_l/1.2/\texttt{len\_l}/blue!5,
      resp_len_h/1.2/\texttt{len\_h}/blue!5,
      resp_payload/2.15/\texttt{payload}/orange!10,
      resp_checksum/1.55/\texttt{checksum}/violet!8
    }{
      \ifnum\i=1
        \node[framebyte,minimum width=\w cm,fill=\color,anchor=west] (\name) at (0,-1.4) {\text};
      \else
        \node[framebyte,minimum width=\w cm,fill=\color,anchor=west] (\name) at (resp_prev.east) {\text};
      \fi
      \coordinate (resp_prev) at (\name.east);
    }
    \draw[decorate,decoration={brace,amplitude=5pt},teal!70!black] (head_a.north west) -- (head_b.north east);
    \node[teal!70!black,font=\scriptsize,above=0.22cm of head_a.north east,anchor=south west] {命令帧头};
    \draw[decorate,decoration={brace,mirror,amplitude=5pt},teal!70!black] (resp_head_a.south west) -- (resp_head_b.south east);
    \node[teal!70!black,font=\scriptsize,below=0.22cm of resp_head_a.south east,anchor=north west] {应答帧头};
  \end{tikzpicture}
  \caption{UART 命令帧和应答帧格式}
  \label{fig:uart_frame}
\end{figure}

\begin{table}[H]
  \centering
  \caption{UART 命令码定义}
  \label{tab:uart_command}
  \begin{tabular}{cccc}
    \toprule
    命令码 & 方向 & 载荷内容 & 功能 \\
    \midrule
    \texttt{0x01} & PC 到 FPGA & 二进制 bit 数据 & 请求 FPGA 编码 \\
    \texttt{0x01} & FPGA 到 PC & HDB3 符号序列 & 返回编码结果 \\
    \texttt{0x02} & PC 到 FPGA & HDB3 符号序列 & 请求 FPGA 译码 \\
    \texttt{0x02} & FPGA 到 PC & 二进制 bit 序列 & 返回译码结果 \\
    \bottomrule
  \end{tabular}
\end{table}

\subsection{DAC 输出方案}

DA0 输出 HDB3 编码波形（双极性 RZ），DA1 输出译码恢复的二进制波形（单极性 NRZ）。FPGA 内部编码到译码存在流水线延迟，但 DAC 模块采用双缓存预填充——两路数据各自完整写入缓存后再从地址 0 同步启动播放，示波器上两路波形码元边界严格对齐。

DAC 播放模块 \texttt{dac\_playback} 内含两个 \texttt{dac\_wave\_ram} 实例（各 $2048 \times 8$ bit，推断为 BSRAM）。采样计数器按 500 个系统时钟为一个码元周期（对应 100 ksym/s），利用 \texttt{rz\_phase} 信号控制前半码元输出有效电平、后半码元归零，实现 DA0 的 RZ 波形；DA1 则整码元保持输出电平。

\begin{figure}[H]
  \centering
  \begin{tikzpicture}[x=0.82cm,y=0.82cm,font=\small]
    \foreach \x in {0,1,...,10}{
      \draw[wavegrid] (\x,-1.5) -- (\x,2.6);
      \node[font=\scriptsize,gray!70] at (\x+0.5,-1.35) {\x};
    }
    \draw[->] (0,-1.15) -- (10.8,-1.15) node[right]{码元时间};

    \node[anchor=east,font=\bfseries] at (-0.82,1.58) {DA0};
    \node[anchor=east,font=\bfseries] at (-0.82,0.28) {DA1};

    \draw[thick,red!70!black]
      (0,1.45) -- (0,1.85) -- (0.5,1.85) -- (0.5,1.45)
      -- (2,1.45) -- (2,1.05) -- (2.5,1.05) -- (2.5,1.45)
      -- (3,1.45) -- (3,1.85) -- (3.5,1.85) -- (3.5,1.45)
      -- (4,1.45) -- (4,1.05) -- (4.5,1.05) -- (4.5,1.45)
      -- (7,1.45) -- (7,1.85) -- (7.5,1.85) -- (7.5,1.45)
      -- (8,1.45) -- (8,1.05) -- (8.5,1.05) -- (8.5,1.45)
      -- (9,1.45) -- (9,1.85) -- (9.5,1.85) -- (9.5,1.45) -- (10,1.45);

    \draw[thick,blue!70!black]
      (0,0.15) -- (1,0.15) -- (1,-0.65) -- (2,-0.65) -- (2,0.15)
      -- (5,0.15) -- (5,-0.65) -- (7,-0.65) -- (7,0.15) -- (10,0.15);

    \node[font=\footnotesize,red!70!black]  at (10.3,1.82) {RZ};
    \node[font=\footnotesize,blue!70!black] at (10.5,0.32) {NRZ};
  \end{tikzpicture}
  \caption{DA0 与 DA1 输出波形关系}
  \label{fig:dac_waveform}
\end{figure}

\section{软硬件实现}

\subsection{HDB3 编码器实现}

FPGA 编码器采用流水线方式，从顶层缓存中逐 bit 读取数据，内部维护四个关键寄存器协同完成编码。

\texttt{zero\_cnt} 为连续零计数器（位宽 2 bit），累计当前已出现的连续零个数；\texttt{ami\_pol} 为 AMI 极性标志，指示下一次普通脉冲应为正（0）还是负（1）；\texttt{pulse\_parity} 为奇偶标志，统计自上一次破坏脉冲 \texttt{V} 以来出现的普通 1 个数的奇偶性（0 为偶数，1 为奇数）；\texttt{last\_pol} 保存最近一次破坏脉冲 \texttt{V} 的极性，用于下一次四连零替换时确定新脉冲的初始极性方向。

编码器内部划分三种工作状态：\texttt{S\_IDLE}（空闲等待）、\texttt{S\_READ}（读取编码）和 \texttt{S\_SUBST}（替换输出）。当所有 bit 处理完成后还会短暂进入 \texttt{S\_FLUSH} 状态以输出末尾尚未成组的零符号。

在 \texttt{S\_READ} 状态下，编码器逐 bit 处理。为提高吞吐率，编码器对连续的零不立即输出，而是将其计入 \texttt{zero\_cnt} 后暂缓——因为零是否直接从输出口送出，取决于它是否属于四连零组。处理规则为：

\begin{itemize}
  \item 输入 bit 为 1 时，若 \texttt{zero\_cnt} 不为零，先进入 \texttt{S\_FLUSH} 状态将暂存零逐一输出，后返回 \texttt{S\_READ} 再处理该 bit 1；若无暂存零，则按 AMI 规则直接输出普通正脉冲 \texttt{+1} 或普通负脉冲 \texttt{-1}，翻转 \texttt{ami\_pol}，	exttt{pulse\_parity} 取反，\texttt{zero\_cnt} 归零。
  \item 输入 bit 为 0 时，\texttt{zero\_cnt} 加 1，当前符号暂不输出。当 \texttt{zero\_cnt} 达到 4 时触发替换；当中途出现 1 时，先将之前暂存的零以符号形式依次输出（\texttt{S\_FLUSH}），再输出该普通脉冲。
\end{itemize}

进入 \texttt{S\_SUBST} 后，编码器根据 \texttt{pulse\_parity} 判断替换形式：若为偶数，采用 \texttt{B00V}，其中 \texttt{B} 按 AMI 规则交替、\texttt{V} 与 \texttt{B} 同极性以破坏交替规律；若为奇数，采用 \texttt{000V}，\texttt{V} 取上一次破坏脉冲的反极性。替换完成后，将四个符号写入 \texttt{subst\_buf} 依次输出，清零 \texttt{zero\_cnt} 和 \texttt{pulse\_parity}，更新 \texttt{last\_pol} 和 \texttt{ami\_pol}，返回 \texttt{S\_READ} 继续处理。

编码结果写入两条通路：一路写入应答缓存 \texttt{resp\_buf}，供 UART 应答模块发回上位机；另一路写入 DA0 波形缓存 \texttt{dac\_buf0}，供 DAC 播放模块驱动 DA0 输出。编码模式下，编码结果还同步送入译码器完成回环译码，译码输出写入 \texttt{dac\_buf1}，使 DA1 能够展示译码恢复波形，与 DA0 编码波形同步对比。

以下为编码器四连零替换核心的 Verilog 实现，对应奇偶判断和符号替换逻辑：

\begin{lstlisting}[style=VerilogStyle,caption={FPGA HDB3 四连零替换核心},label={lst:encoder_core}]
// 检测连续四个 0 并生成替换符号
if (zero_cnt == 2'd3) begin
    if (pulse_parity == 1'b0) begin
        // 偶数普通脉冲采用 B00V 替换
        if (last_pol == 1'b0) begin
            subst_buf[0] <= SYM_NB;
            subst_buf[3] <= SYM_NV;
            ami_pol      <= 1'b0;
            last_pol     <= 1'b1;
        end
        else begin
            subst_buf[0] <= SYM_PB;
            subst_buf[3] <= SYM_PV;
            ami_pol      <= 1'b1;
            last_pol     <= 1'b0;
        end
    end
    else begin
        // 奇数普通脉冲采用 000V 替换
        subst_buf[0] <= SYM_0;
        subst_buf[3] <= (last_pol == 1'b0) ? SYM_PV : SYM_NV;
        ami_pol      <= ~last_pol;
    end
    subst_buf[1] <= SYM_0;
    subst_buf[2] <= SYM_0;
    subst_idx    <= 2'd0;
    zero_cnt     <= 2'd0;
    pulse_parity <= 1'b0;
    state        <= S_SUBST;
end
\end{lstlisting}

\subsection{HDB3 译码器实现}

译码器采用直接符号映射方式，按符号逐项处理输入 HDB3 序列，不反向查找 \texttt{B00V} 或 \texttt{000V} 的组合位置。其译码规则严格基于表 \ref{tab:symbol_definition} 中的符号定义：凡来源于原始输入 bit 1 的符号（普通正脉冲 \texttt{+1} 和普通负脉冲 \texttt{-1}）译码为 1；凡来源于原始输入 bit 0 的符号（零电平 \texttt{0}、破坏脉冲 \texttt{+V/-V}、补信脉冲 \texttt{+B/-B}）译码为 0。

这种直译方式与编码器的奇偶法相互独立。编码器按照内部奇偶状态决定 \texttt{B00V} 还是 \texttt{000V}，译码器不关心符号属于哪种替换模式，只关心该符号对应原始 bit 是 0 还是 1。编码器替换策略的改变不会影响译码结果，两个模块在逻辑上完全解耦，这对 FPGA 的分模块开发和独立验证非常有利。

在译码模式下，上位机将 HDB3 符号序列通过 UART 发送给 FPGA，顶层主控将输入符号写入 DA0 缓存（显示输入符号的波形），同时将译码输出的 bit 序列写入 DA1 缓存（显示译码恢复的二进制波形）和应答缓存。两路波形码元边界对齐，便于对比观察符号输入和 bit 恢复的对应关系。

\subsection{主控状态机实现}

顶层状态机负责管理一次命令从接收到应答的完整生命周期。当 UART 接收模块完成 \texttt{55 AA} 命令帧解析后，状态机从 \texttt{M\_IDLE} 跳转到编码或译码路径，根据命令码依次完成数据处理、DAC 缓存写入、DAC 播放启动和 UART 应答发送。各状态功能如下：

\textbf{空闲（M\_IDLE）}：等待 \texttt{packet\_parser} 完成命令帧解析。解析完成后，根据命令码 \texttt{0x01}（编码）或 \texttt{0x02}（译码）跳转到对应准备状态；若帧解析失败或参数非法，跳转到错误应答状态。

\textbf{编码准备（M\_ENC\_PREP）}：将载荷中的 bit 数据搬入编码器输入缓存，同时复制到软件参考区供后续比对。

\textbf{编码运行（M\_ENC\_RUN）}：启动编码器，等待编码器逐 bit 完成全部处理，编码结果同步写入应答缓存和 DA0 缓存。

\textbf{回环译码（M\_ENC\_LOOP）}：仅在编码模式下存在。将编码器输出的 HDB3 符号送入译码器进行回环译码，译码输出写入 DA1 缓存，使示波器可同时观察编码波形和译码恢复波形。

\textbf{译码准备（M\_DEC\_PREP）}：将载荷中的 HDB3 符号数据搬入译码器输入缓存和 DA0 缓存。

\textbf{译码运行（M\_DEC\_RUN）}：启动译码器，等待译码器逐符号完成全部处理，译码结果写入应答缓存和 DA1 缓存。

\textbf{启动 DAC（M\_START\_DAC）}：将 DA0 和 DA1 波形缓存的数据长度写入 DAC 播放模块，触发循环播放。DAC 播放独立运行，状态机不等待播放完成。

\textbf{写应答（M\_RESP\_WRITE）}：将结果按 \texttt{AA 55} 应答帧格式写入 \texttt{resp\_buf}，包含帧头、命令码、状态码、载荷长度、载荷数据和校验和。

\textbf{发送应答（M\_RESP\_SEND）}：触发 UART 发送模块逐字节发送 \texttt{resp\_buf}。发送完成后回到 \texttt{M\_IDLE}。

\textbf{错误应答（M\_ERR\_SEND）}：生成仅含帧头和错误状态码的简短应答帧并发送，发送完成后回到 \texttt{M\_IDLE}。

编码和译码路径在各自运行完毕后汇合到 \texttt{M\_START\_DAC}，之后共享 DAC 启动、应答写入、应答发送三条公共路径。错误路径从 \texttt{M\_IDLE} 直接进入 \texttt{M\_ERR\_SEND}，不经过任何数据处理和 DAC 环节。状态机流程如图 \ref{fig:main_fsm} 所示。

\begin{figure}[H]
  \centering
  \resizebox{0.94\textwidth}{!}{%
  \begin{tikzpicture}[font=\small]
    \node[smallblock] (idle) at (0,0) {空闲\\\texttt{M\_IDLE}};
    \node[smallblock] (encprep) at (2.9,1.00) {编码准备\\\texttt{M\_ENC\_PREP}};
    \node[smallblock] (encrun) at (5.5,1.00) {编码运行\\\texttt{M\_ENC\_RUN}};
    \node[smallblock] (encloop) at (8.1,1.00) {回环译码\\\texttt{M\_ENC\_LOOP}};
    \node[smallblock] (decprep) at (2.9,-0.80) {译码准备\\\texttt{M\_DEC\_PREP}};
    \node[smallblock] (decrun) at (5.5,-0.80) {译码运行\\\texttt{M\_DEC\_RUN}};
    \coordinate (merge) at (10.0,0.10);
    \node[smallblock] (dac) at (11.8,0.10) {启动 DAC\\\texttt{M\_START\_DAC}};
    \node[smallblock] (write) at (14.4,0.10) {写应答\\\texttt{M\_RESP\_WRITE}};
    \node[smallblock] (send) at (16.7,0.10) {发送应答\\\texttt{M\_RESP\_SEND}};
    \node[smallblock] (err) at (9.6,-2.35) {错误应答\\\texttt{M\_ERR\_SEND}};

    \coordinate (fork) at (1.35,0);
    \draw[arrow] (idle.east) -- (fork) |- (encprep.west);
    \draw[arrow] (encprep) -- (encrun);
    \draw[arrow] (encrun) -- (encloop);
    \draw[arrow] (encloop.east) -- ++(0.45,0) |- (merge);

    \draw[arrow] (fork) |- (decprep.west);
    \draw[arrow] (decprep) -- (decrun);
    \draw[arrow] (decrun.east) -- ++(2.00,0) |- (merge);

    \fill[black!70] (merge) circle (1.2pt);
    \draw[arrow] (merge) -- (dac.west);
    \draw[arrow] (dac.east) -- (write.west);
    \draw[arrow] (write.east) -- (send.west);

    \draw[arrow] (send.north) -- ++(0,1.5) -- ++(-16.2,0) |- (idle.north);

    \draw[arrow] (idle.south) -- ++(0,-2.30) -- node[font=\small,above,pos=0.45] {帧错误 / 参数错误} (err.west);
    \draw[arrow] (err.east) -- ++(5.20,0) -| (send.south);
  \end{tikzpicture}
  }
  \caption{顶层主控状态机流程}
  \label{fig:main_fsm}
\end{figure}

\subsection{DAC 播放实现}

DAC 播放模块 \texttt{dac\_playback} 负责将波形缓存中的符号序列转换为模拟电平输出。其核心逻辑包括 HDB3 符号到 DAC 电平的映射，以及 RZ/NRZ 两种输出格式的时序控制：

\begin{lstlisting}[style=VerilogStyle,caption={DAC 码元播放与电平映射核心},label={lst:dac_core}]
// 定义 DAC 输出电平
localparam [7:0] DAC_ZERO = 8'h80;
localparam [7:0] DAC_POS  = 8'h00;
localparam [7:0] DAC_NEG  = 8'hFF;
// 250 个系统时钟翻转一次相位，500 个时钟对应一个码元
wire half_tick = (sample_cnt == 8'd249);
// 将 HDB3 符号映射到 DA0 电平
function [7:0] hdb3_symbol_to_dac;
    input [7:0] sym;
    begin
        case (sym)
            8'h01, 8'h03, 8'h05: hdb3_symbol_to_dac = DAC_POS;
            8'h02, 8'h04, 8'h06: hdb3_symbol_to_dac = DAC_NEG;
            default:             hdb3_symbol_to_dac = DAC_ZERO;
        endcase
    end
endfunction
// DA0 输出 HDB3 RZ，DA1 输出译码 NRZ
always @(posedge clk or negedge rst_n) begin
    if (!rst_n) begin
        DA0_Data <= DAC_ZERO;
        DA1_Data <= DAC_ZERO;
    end
    else if (sample_cnt == 8'd0) begin
        DA0_Data <= rz_phase ? DAC_ZERO : hdb3_symbol_to_dac(rd_data0);
        DA1_Data <= (rd_data1 == 8'h00) ? DAC_ZERO : DAC_POS;
    end
end
\end{lstlisting}

其中 \texttt{hdb3\_symbol\_to\_dac} 函数将 6 种 HDB3 符号统一映射到正、负、零三种模拟电平：正符号（\texttt{+1}、\texttt{+V}、\texttt{+B}）输出 \texttt{8'h00}，负符号（\texttt{-1}、\texttt{-V}、\texttt{-B}）输出 \texttt{8'hFF}，零符号输出 \texttt{8'h80}。DA0 通过 \texttt{rz\_phase} 实现半码元归零，DA1 则整码元保持 \texttt{rd\_data1} 的非零判断结果。


\subsection{上位机实现}

上位机采用 C\# WPF（.NET 6.0）框架编写，使用 MVVM 模式组织代码。主界面分为编码测试区和译码测试区两个 Tab：编码区输入二进制序列（仅接受 0 和 1），显示软件期望的 HDB3 符号结果和 FPGA 实际返回结果；译码区输入 HDB3 符号序列，显示软件期望的 bit 结果和 FPGA 实际返回结果。比对结果通过颜色和文本状态标识：一致显示绿色，不一致显示红色。

软件核心由三个类协同工作：

\texttt{HDB3Codec.cs} 实现 HDB3 软件参考编译码算法，与 FPGA 侧的 Verilog 实现采用相同的奇偶法规则。编码器与 FPGA 端类似，通过 3 符号延迟缓冲暂存零，遇到连续四个零时依据 \texttt{B00V/000V} 替换规则生成 HDB3 符号；译码器按符号类型进行直译映射。该类还提供 bit 序列的字节打包（\texttt{PackBits}）和符号与字符串的相互转换（\texttt{SymbolToString}）等工具方法。

\texttt{SerialService.cs} 基于 \texttt{System.IO.Ports.SerialPort} 实现串口管理。帧解析采用字节流状态机：在接收缓冲中扫描 \texttt{AA 55} 帧头，定位后顺序读取命令码、状态码、载荷长度和载荷数据，校验和验证通过后通过事件回调通知 ViewModel 层。\texttt{DataReceived} 事件在后台线程中触发，界面更新通过 \texttt{Dispatcher.Invoke} 调度到 UI 线程。

\texttt{MainViewModel.cs} 负责界面数据绑定和命令调度。编码时先将输入字符串转为 bit 数组，调用 \texttt{Encode} 生成软件期望，再发送编码命令（\texttt{0x01}）。收到应答后逐符号比对，结果以颜色标记正确或错误并绑定到界面列表。译码流程同理（\texttt{0x02}）。界面还支持编码结果一键导入译码区和日志输出等功能。

以上三个类协同完成一次编码命令的流程如图 \ref{fig:pc_flow} 所示。

\begin{figure}[H]
  \centering
  \resizebox{0.94\textwidth}{!}{%
  \begin{tikzpicture}[font=\footnotesize,node distance=0.55cm]
    \node[smallblock,minimum width=1.6cm] (in) {输入序列};
    \node[smallblock,right=of in,minimum width=1.6cm] (sw) {软件编码};
    \node[smallblock,right=of sw,minimum width=1.6cm] (tx) {发送命令};
    \node[smallblock,right=of tx,minimum width=1.6cm] (fpga) {FPGA};
    \node[smallblock,right=of fpga,minimum width=1.6cm] (rx) {接收应答};
    \node[smallblock,right=of rx,minimum width=1.6cm] (cmp) {逐符号比对};
    \node[smallblock,right=of cmp,minimum width=1.6cm] (ui) {界面显示};

    \draw[arrow] (in) -- (sw);
    \draw[arrow] (sw) -- (tx);
    \draw[arrow] (tx) -- (fpga);
    \draw[arrow] (fpga) -- (rx);
    \draw[arrow] (rx) -- (cmp);
    \draw[arrow] (cmp) -- (ui);
  \end{tikzpicture}
  }
  \caption{上位机编码命令处理流程}
  \label{fig:pc_flow}
\end{figure}

其中，\texttt{MainViewModel} 发送命令的核心 C\# 逻辑如下，体现了软件期望生成、载荷打包和命令发送三个关键步骤：

\begin{lstlisting}[style=CSharpStyle,caption={上位机编码命令发送核心},label={lst:pc_send_core}]
// 提取输入 bit 并生成软件期望结果
var chars = EncInput.Where(c => c == '0' || c == '1').ToArray();
var bits = chars.Select(c => c == '1').ToArray();
byte[] expected = HDB3Codec.Encode(bits);

// 按协议构建编码命令载荷
int bitCnt = bits.Length;
byte[] packed = HDB3Codec.PackBits(bits);
var payload = new byte[2 + packed.Length];
payload[0] = (byte)(bitCnt & 0xFF);
payload[1] = (byte)((bitCnt >> 8) & 0xFF);
Array.Copy(packed, 0, payload, 2, packed.Length);

// 发送 0x01 编码命令并刷新界面期望结果
_serial.SendCommand(0x01, payload);
EncExpected.Clear();
foreach (var symbol in expected)
{
    EncExpected.Add(new ObservableSymbol {
        Value = HDB3Codec.SymbolToString(symbol),
        IsCorrect = true
    });
}
\end{lstlisting}

FPGA 应答返回后，比对结果通过界面列表展示，正确项与错误项以不同颜色区分。上位机运行界面如图 \ref{fig:pc_interface} 所示。

\begin{figure}[H]
  \centering
  \includegraphics[width=0.95\textwidth]{上位机界面.png}
  \caption{上位机测试界面}
  \label{fig:pc_interface}
\end{figure}

\newpage
\section{仿真与实测验证}

\subsection{独立仿真工程验证}

为在纯逻辑层面验证 HDB3 编译码算法的正确性，本文在 Vivado 中搭建了独立仿真工程，得到两类结果：一是 Vivado 自动生成的 RTL 结构框图（图 \ref{fig:rtl_sim}），二是 ModelSim 仿真波形（图 \ref{fig:modelsim_wave}）。

RTL 结构框图由 Vivado 综合后自动绘制，展示了顶层模块 \texttt{hdb3\_top} 的模块连接关系和端口信号。从图中可以直观看到编码器、译码器和测试激励模块 \texttt{tb\_hdb3} 之间的信号流向：\texttt{clk} 和 \texttt{rst\_n} 接入各模块，\texttt{data\_in} 送入编码器，编码器输出的 HDB3 双线信号经顶层处理后分别以 \texttt{hdb3\_pos}/\texttt{hdb3\_neg} 形式引出，同时送入译码器得到 \texttt{data\_out} 和 \texttt{data\_valid}。需要说明的是，该 RTL 图反映的是纯仿真环境的模块结构，不包含 UART 通信、DAC 播放等最终系统的外围模块，因此本文仅将其作为仿真补充，不用于说明最终 FPGA 系统划分。

\begin{figure}[H]
  \centering
  \includegraphics[width=1.0\textwidth]{仿真的RTL框图.png}
  \caption{独立仿真工程 RTL 结构图}
  \label{fig:rtl_sim}
\end{figure}

仿真输入序列为 \texttt{1000010000111100}，其中包含四连零、普通 1 和连续 1，便于同时观察 HDB3 替换、极性交替和译码恢复。图 \ref{fig:modelsim_wave} 为 ModelSim 仿真波形，各信号端口含义如表 \ref{tab:modelsim_ports} 所示。

\begin{figure}[H]
  \centering
  \includegraphics[width=1.0\textwidth]{Modelsim仿真图.png}
  \caption{ModelSim HDB3 编译码波形仿真}
  \label{fig:modelsim_wave}
\end{figure}

仿真时钟 \texttt{clk} 按半码元周期翻转，因此每两个 \texttt{clk} 周期对应一个完整码元。各信号端口含义如表 \ref{tab:modelsim_ports} 所示。

\begin{table}[H]
  \centering
  \caption{ModelSim 仿真端口说明}
  \label{tab:modelsim_ports}
  \begin{tabularx}{\textwidth}{>{\ttfamily}c c X}
    \toprule
    端口 & 方向 & 含义 \\
    \midrule
    clk & 输入 & 半码元仿真时钟，testbench 中周期为 10 ns \\
    rst\_n & 输入 & 低有效复位信号 \\
    data\_in & 输入 & testbench 输入的单极性二进制序列 \\
    hdb3\_sym\_pos / hdb3\_sym\_neg & 内部观测 & 编码器输出的完整码元极性，送入译码器作为 HDB3 符号输入 \\
    hdb3\_valid & 内部观测 & HDB3 符号有效信号，用于限定译码器采样时刻 \\
    hdb3\_pos / hdb3\_neg & 输出观测 & 顶层根据\texttt{phase}门控后的RZ脉冲 \\
    data\_out & 输出 & 译码器恢复出的二进制 bit \\
    data\_valid & 输出 & 译码输出有效信号，testbench 据此进行逐 bit 比对 \\
    code\_error & 输出 & HDB3 双线表示非法时置位 \\
    \bottomrule
  \end{tabularx}
\end{table}

从波形结果可以看出，输入序列中的四连零位置会在 HDB3 符号极性波形上产生替换脉冲；归零输出只在码元前半段保持有效，后半段回到零电平；经过译码器流水线延迟后，\texttt{data\_out} 能够按 \texttt{data\_valid} 节奏恢复原始 bit 序列。最终系统正确性仍需结合上位机比对和示波器实测进行判断。

\subsection{编码功能实测}

编码功能测试中，上位机输入二进制序列，先由软件参考算法生成期望 HDB3 符号序列，再通过 UART 将输入发送给 FPGA。FPGA 返回编码结果后，上位机逐符号比较二者是否一致。与此同时，DA0 输出 HDB3 编码波形，DA1 输出回环译码波形。

\begin{table}[H]
  \centering
  \caption{编码测试用例}
  \label{tab:encode_cases}
  \begin{tabularx}{\textwidth}{c c X}
    \toprule
    测试项 & 输入序列 & 验证重点 \\
    \midrule
    全 0 测试 & \texttt{0000000000000000} & 连续四零替换和译码恢复为 0 \\
    全 1 测试 & \texttt{1111111111111111} & AMI 正负脉冲交替 \\
    10 交替测试 & \texttt{1010101010101010} & 普通 1 的交替脉冲和 0 的间隔 \\
    典型序列测试 & \texttt{1000010000111100} & 普通脉冲、四连零替换和连续 1 综合验证 \\
    长 0 测试 & \texttt{1100000000110000} & 多组四连零替换 \\
    编码随机测试 & \texttt{10111000010000001} & 一般序列下的软件与 FPGA 一致性 \\
    \bottomrule
  \end{tabularx}
\end{table}

\subsubsection{全 0 序列测试}

全 0 序列(0000000000000000)中没有普通 1，如果采用普通 AMI 码会长时间保持零电平。HDB3 编码会周期性插入替换脉冲，从而避免长连零。图 \ref{fig:all_zero_test} 中，上位机显示 FPGA 返回结果与软件参考结果一致；示波器上 DA0 出现替换脉冲，DA1 保持低电平，说明译码结果恢复为 0。

\begin{figure}[H]
  \centering
  \begin{subfigure}{0.50\textwidth}
    \centering
    \includegraphics[width=\linewidth]{全0测试-上位机.png}
    \caption{上位机比对结果}
  \end{subfigure}
  \hfill
  \begin{subfigure}{0.48\textwidth}
    \centering
    \includegraphics[width=\linewidth]{全0测试-FPGA.jpg}
    \caption{示波器波形}
  \end{subfigure}
  \caption{全 0 序列测试结果}
  \label{fig:all_zero_test}
\end{figure}

\subsubsection{全 1 序列测试}

全 1 序列(1111111111111111)不触发四连零替换，主要验证 AMI 极性交替规则。图 \ref{fig:all_one_test} 中，FPGA 返回结果与软件参考一致；示波器上 DA0 呈现连续正负交替脉冲，DA1 保持高电平，符合全 1 序列译码结果。

\begin{figure}[H]
  \centering
  \begin{subfigure}{0.50\textwidth}
    \centering
    \includegraphics[width=\linewidth]{全1测试-上位机.png}
    \caption{上位机比对结果}
  \end{subfigure}
  \hfill
  \begin{subfigure}{0.48\textwidth}
    \centering
    \includegraphics[width=\linewidth]{全1测试-FPGA.jpg}
    \caption{示波器波形}
  \end{subfigure}
  \caption{全 1 序列测试结果}
  \label{fig:all_one_test}
\end{figure}

\newpage
\subsubsection{10 交替序列测试}

10 交替序列(1010101010101010)用于验证普通 1 的正负交替和 0 的零电平间隔。该序列不产生四连零替换，因此 DA0 中应当表现为间隔出现的正负脉冲，DA1 中应当表现为 1、0 交替的 NRZ 波形。测试结果如图 \ref{fig:alternate_test} 所示。

\begin{figure}[H]
  \centering
  \begin{subfigure}{0.49\textwidth}
    \centering
    \includegraphics[width=\linewidth]{10交替测试-上位机.png}
    \caption{上位机比对结果}
  \end{subfigure}
  \hfill
  \begin{subfigure}{0.48\textwidth}
    \centering
    \includegraphics[width=\linewidth]{10交替测试-FPGA.jpg}
    \caption{示波器波形}
  \end{subfigure}
  \caption{10 交替序列测试结果}
  \label{fig:alternate_test}
\end{figure}

\subsubsection{典型混合序列测试}

典型混合序列(1000010000111100)同时包含普通 1、连续 0 和连续 1，是最能综合体现 HDB3 编码规则的测试用例。图 \ref{fig:mixed_test} 中，上位机比对结果一致；示波器波形中既能观察到 AMI 正负交替，也能观察到四连零替换产生的脉冲，DA1 能够恢复原始输入节奏。

\begin{figure}[H]
  \centering
  \begin{subfigure}{0.49\textwidth}
    \centering
    \includegraphics[width=\linewidth]{典型序列测试-上位机.png}
    \caption{上位机比对结果}
  \end{subfigure}
  \hfill
  \begin{subfigure}{0.48\textwidth}
    \centering
    \includegraphics[width=\linewidth]{典型序列测试-FPGA.jpg}
    \caption{示波器波形}
  \end{subfigure}
  \caption{典型混合序列测试结果}
  \label{fig:mixed_test}
\end{figure}

\newpage
\subsubsection{长 0 序列测试}

长 0 序列(1100000000110000)用于验证多组连续四零替换。图 \ref{fig:long_zero_test} 中，FPGA 返回结果与软件参考一致；示波器上 DA0 能看到多组替换脉冲，说明 HDB3 编码有效抑制了长连零。

\begin{figure}[H]
  \centering
  \begin{subfigure}{0.49\textwidth}
    \centering
    \includegraphics[width=\linewidth]{长0序列测试-上位机.png}
    \caption{上位机比对结果}
  \end{subfigure}
  \hfill
  \begin{subfigure}{0.48\textwidth}
    \centering
    \includegraphics[width=\linewidth]{长0序列测试-FPGA.jpg}
    \caption{示波器波形}
  \end{subfigure}
  \caption{长 0 序列测试结果}
  \label{fig:long_zero_test}
\end{figure}

\subsubsection{编码随机序列测试}

随机序列测试用于验证系统在一般输入条件下的稳定性。本次实测输入序列为 \texttt{10111000010000001}。图 \ref{fig:encode_random_test} 中，上位机软件结果和 FPGA 返回结果一致，示波器波形也能与输入序列节奏对应，说明编码链路、UART 返回和 DAC 输出均工作正常。

\begin{figure}[H]
  \centering
  \begin{subfigure}{0.48\textwidth}
    \centering
    \includegraphics[width=\linewidth]{编码随机测试-上位机.png}
    \caption{上位机比对结果}
  \end{subfigure}
  \hfill
  \begin{subfigure}{0.48\textwidth}
    \centering
    \includegraphics[width=\linewidth]{编码随机测试-FPGA.jpg}
    \caption{示波器波形}
  \end{subfigure}
  \caption{编码随机序列测试结果}
  \label{fig:encode_random_test}
\end{figure}

\newpage
\subsection{译码功能实测}

译码测试中，上位机输入 HDB3 符号序列，软件参考算法先得到期望 bit 序列，然后 FPGA 对相同符号序列进行译码并返回 bit 结果。上位机逐 bit 比对软件期望和 FPGA 返回值。图 \ref{fig:decode_import} 展示了由编码结果导入译码区的操作方式。

\begin{figure}[H]
  \centering
  \includegraphics[width=0.90\textwidth]{译码可直接导入-以全1为例.png}
  \caption{编码结果导入译码区示例}
  \label{fig:decode_import}
\end{figure}

随机译码测试采用的实际输入和期望输出如表 \ref{tab:decode_random_case} 所示。该序列包含普通正负脉冲、补信脉冲和破坏脉冲，可用于检查译码器对不同 HDB3 符号的映射是否一致。

\begin{table}[H]
  \centering
  \caption{随机译码实测序列}
  \label{tab:decode_random_case}
  \begin{tabularx}{\textwidth}{c>{\ttfamily\raggedright\arraybackslash}X}
    \toprule
    项目 & 序列 \\
    \midrule
    输入 HDB3 符号 & +1 0 0 -1 +B 0 0 +V -1 0 0 0 -V 0 0 +1 -1 0 \\
    期望 bit 输出 & 100100001000000110 \\
    \bottomrule
  \end{tabularx}
\end{table}

随机译码测试结果如图 \ref{fig:decode_random_test} 所示。上位机结果表明 FPGA 译码输出与软件参考一致；示波器中 DA0 显示输入 HDB3 符号波形，DA1 显示译码后的 NRZ bit 波形，两路波形在码元边界上保持同步。

\begin{figure}[H]
  \centering
  \begin{subfigure}{0.48\textwidth}
    \centering
    \includegraphics[width=\linewidth]{解码随机测试-上位机.png}
    \caption{上位机比对结果}
  \end{subfigure}
  \hfill
  \begin{subfigure}{0.48\textwidth}
    \centering
    \includegraphics[width=\linewidth]{解码随机测试-FPGA.jpg}
    \caption{示波器波形}
  \end{subfigure}
  \caption{解码随机序列测试结果}
  \label{fig:decode_random_test}
\end{figure}

\subsection{测试结果汇总}

各组测试结果如表 \ref{tab:test_summary} 所示。可以看出，不同类型输入下，上位机软件结果与 FPGA 返回结果均保持一致，示波器波形也符合对应测试工况的预期。

\begin{table}[H]
  \centering
  \caption{测试结果汇总}
  \label{tab:test_summary}
  \begin{tabularx}{\textwidth}{cX X c}
    \toprule
    测试项目 & 上位机比对 & 示波器观察 & 结论 \\
    \midrule
    全 0 编码 & 软件结果与 FPGA 返回一致 & DA0替换脉冲，DA1为低 & 通过 \\
    全 1 编码 & 软件结果与 FPGA 返回一致 & DA0正负交替，DA1为高 & 通过 \\
    10 交替编码 & 软件结果与 FPGA 返回一致 & DA0间隔脉冲，DA1交替 & 通过 \\
    典型序列编码 & 软件结果与 FPGA 返回一致 & 普通和替换脉冲均可观察 & 通过 \\
    长 0 编码 & 软件结果与 FPGA 返回一致 & 多组替换脉冲明显 & 通过 \\
    编码随机测试 & 软件结果与 FPGA 返回一致 & 波形节奏与输入相符 & 通过 \\
    解码随机测试 & 软件结果与 FPGA 返回一致 & DA1 恢复译码 bit & 通过 \\
    \bottomrule
  \end{tabularx}
\end{table}

本实验没有搭建随机噪声信道，因此不绘制传统误码率曲线。验证方式采用固定序列和随机序列的逐 bit、逐符号比对。在当前测试范围内，软件参考结果与 FPGA 返回结果一致，实测误码数为 0。

\newpage
\section{问题分析与改进}

\subsection{仿真工程与最终系统的区别}

独立仿真工程基于 Vivado/ModelSim 环境构建，顶层模块直接接收 testbench 提供的激励信号，编码器和译码器以串联流水线方式工作。该工程的核心目的是在纯逻辑层面验证 HDB3 编译码算法的正确性，不包含任何 I/O 接口和通信协议。

最终 Gowin FPGA 工程以 \texttt{hdb3\_top} 为顶层，在编译码核心逻辑外围增加了完整的通信和输出子系统：UART 字节收发模块（\texttt{uart\_byte\_rx/uart\_byte\_tx}）负责与上位机通信，命令帧解析模块（\texttt{packet\_parser}）负责提取载荷和控制字，应答发送模块（\texttt{response\_tx}）负责封装 \texttt{AA 55} 应答帧，DAC 播放模块（\texttt{dac\_playback}）负责将符号序列转换为模拟波形。仿真工程中的 \texttt{hdb3\_top} 与最终工程中的 \texttt{hdb3\_top} 并非同一模块——仿真顶层的功能范围仅相当于最终工程中 \texttt{hdb3\_encoder} 与 \texttt{hdb3\_decoder} 的组合。

因此，本文将仿真 RTL 图和 ModelSim 波形放在仿真验证部分，作为编译码逻辑正确性的补充说明；而图 \ref{fig:fpga_structure} 展示的最终 FPGA 工程模块关系图才是实际系统的真实结构，读者不应将二者混淆。

\subsection{四连零替换规则的统一}

HDB3 规则在不同教材中描述方式差异较大。常见的一种叙述是：首次遇到四连零时固定采用 \texttt{000V}，V的极性与前一个1的极性相同，后续则根据两次替换组之间的 1 个数奇偶性判断；另一种叙述（本设计采用的奇偶法）则是统一按自上一次破坏脉冲以来普通 1 的奇偶个数判断，初始奇偶状态为偶数。

两种规定在多数序列下输出相同，但在全零序列的首个替换组上存在差异：按"首次固定 \texttt{000V}"的规定，全零序列的首个替换组为 \texttt{000V}，首组之后若两个普通 1 的奇偶计数为偶数则使用 \texttt{B00V}；而按奇偶法，初始状态为偶数，全零序列的首个替换组直接采用 \texttt{B00V}。这是教材规定与产业工程实现之间的一个已知变体，不属于实现错误。

本设计选择奇偶法的原因在于：其判断逻辑与 AMI 极性交替共用同一计数器，在 FPGA 硬件中无需额外的"首次标志"寄存器，实现更加规整。为确保上位机参考结果与 FPGA 返回结果的可比性，上位机 \texttt{HDB3Codec.cs} 和 FPGA \texttt{hdb3\_encoder.v} 均以同一套奇偶法实现，所有测试用例的比对结果均基于统一规则，保证了验证的有效性。


\subsection{FPGA 综合与布局布线速度优化}

在开发过程中，FPGA EDA 工具的综合与布局布线耗时一度过长。分析发现，问题根源在于 \texttt{dac\_wave\_ram.v} 中波形缓存的 RTL 写法。

该模块使用 \texttt{reg [7:0] mem [0:2047]} 声明一个 $2048 \times 8$ bit 的存储阵列，设计意图是让综合工具将其推断为 BSRAM（块状 RAM）。然而，读端口最初写为：

\begin{lstlisting}[style=VerilogStyle,caption={优化前的读端口时序（关键问题）},label={lst:ram_bad}]
always @(posedge clk or negedge rst_n) begin
    if (!rst_n)
        rd_data <= 8'h80;
    else
        rd_data <= mem[rd_addr];
end
\end{lstlisting}

\texttt{always @(posedge clk or negedge rst\_n)} 将异步复位信号引入读数据寄存器。Gowin BSRAM 的 DO 寄存器不支持异步复位，导致综合工具无法将其推断为块 RAM，无奈回退为数千个触发器（$2048 \times 8 = 16384$ 个 FF）实现的寄存器阵列。这不仅大量消耗 LUT 资源，更令综合网表规模膨胀，布局布线的等效门数远超实际需求，编译时间显著变长。

优化方案是将读端口改为纯同步时序：

\begin{lstlisting}[style=VerilogStyle,caption={优化后的读端口（纯同步，可推断为 BSRAM）},label={lst:ram_good}]
always @(posedge clk) begin
    rd_data <= mem[rd_addr];
end
\end{lstlisting}

修改后，综合工具成功将波形缓存数组推断为 BSRAM，LUT 利用率大幅下降，综合与布局布线时间恢复至合理范围。这一经验说明：FPGA 中存储阵列的 RTL 必须遵循目标器件 IP 的推断模板——异步复位、初始值赋值等常见写法都可能阻碍原语推断，导致面积和编译时间的成倍恶化。

复位后读数据初始值未知。但 DAC 播放模块在上层通过播放标志和等待计数器控制输出时序，在数据有效前将两路输出强制置为中点电平，因此读端初始不定态不影响实际输出。

此外，\texttt{hdb3\_top.v} 中的 \texttt{payload\_buf[0:63]} 和 \texttt{sym\_buf[0:63]}（各 $64 \times 8$ bit）规模有限，即使推断为分布式 RAM 也不会造成明显瓶颈，但仍采用同步时序写法以保证代码风格的一致性。

\newpage
\section{总结}

本文完成了一套基于 FPGA 的 HDB3 编译码器系统。系统以 ACX720 开发板（GW5AT-138B）为硬件平台，实现了二进制序列到 HDB3 符号序列的编码、HDB3 符号序列到二进制序列的译码、单 UART 命令/应答通信和双通道 DAC 波形输出。上位机采用 C\# WPF 编写，通过独立的软件参考算法生成期望结果，与 FPGA 返回值进行逐符号比对，形成"软件参考—FPGA 返回—示波器波形"三层验证闭环。

在实现层面，编码器采用奇偶法判断四连零替换形式，以 \texttt{S\_READ}、\texttt{S\_FLUSH}、\texttt{S\_SUBST} 三状态流水线完成零符号延迟和批量替换；译码器采用直接符号映射，与编码器在逻辑上解耦。DAC 播放模块通过双缓存预填充机制消除了编译码流水线延迟，使两路波形在示波器上码元对齐。开发过程中曾因 \texttt{dac\_wave\_ram} 的异步复位写法阻碍 BSRAM 推断，导致综合速度严重劣化，改为纯同步读端口后恢复正常，这一经验对 FPGA 存储阵列的 RTL 设计具有参考意义。

实验结果表明，在全 0、全 1、10 交替、典型混合、长 0、编码随机和解码随机等工况下，上位机软件参考结果与 FPGA 返回结果一致，示波器能够观察到 HDB3 码的正负交替脉冲、四连零替换脉冲和译码恢复波形，FPGA 编译码逻辑、UART 通信和 DAC 输出均达到设计要求。

后续可进一步完善系统级 testbench 覆盖，增加更长时间随机序列测试和异常帧处理测试，并探索将编译码参数（码元速率、序列长度）设计为可动态配置，提升系统的通用性。

\end{document}
